\section{The Manifesto LLM Replication}
\label{sec:v8-manifesto-results}

We instantiate the framework on the Manifesto LLM benchmark of
\citet{BenoitEtAl2025}: \(235\) party platforms scored against
expert-survey means on six seven-point policy dimensions (economic
policy, social policy, immigration, European integration,
environmental policy, decentralization). The source documents are
party election programs from the Comparative Manifestos Project
(CMP, hereafter the Manifesto Project). Before the methodology, we
walk through what a tree state has to carry on this task.

\subsection{What the State Has to Carry}
\label{sec:v8-manifesto-mechanism}

Take the economic-policy scale as a running example. A trusted
coder reads the full party platform and places it on the
seven-point scale; that placement is the oracle \(f^\ast\) we want
a learned tree to recover. To be useful for the same target, a
C-Tree state for a span has to carry whatever evidence could move
that placement: tax promises, spending commitments, privatization
claims, labor-market policy, welfare language, fiscal constraints,
and the qualifications about who benefits or pays.

A scalar score per leaf is too small a state for this work. Suppose
one section of a platform promises a tax cut for small businesses,
and a later section expands health spending. Suppose a third
section makes both conditional on private delivery and on a deficit
rule. The economic position of the platform turns on how those
three pieces fit together. A parent that only inherits two scalar
child scores cannot recover the qualification. A parent that
inherits two policy-evidence states can carry the tax position, the
spending position, and the condition that ties them together.

A C-Tree state is the working memory the tree keeps about a span.
The query that comes out at the end is a number, but the state has
to carry more than a number. In economic policy, a state may record
that a party supports tax cuts for small businesses, protects health
spending, and funds the package through privatization. In
environmental policy, it may carry both emissions targets and
exceptions for energy security. In decentralization, it may need to
distinguish administrative devolution from fiscal autonomy, regional
veto power, and country-specific institutional language. The score
is read off the root, so the state has to survive every merge before
that read is meaningful.

The three local laws name the failure modes. C1 (sufficiency) fails
when a leaf summary drops the evidence the rubric needs. C2
(idempotence) fails when a state drifts after a re-summarization
pass, for example when a conditional tax cut becomes unconditional.
C3 (merge consistency) fails when two child states are each fine on
their own and lose their interaction once merged. C3 is the
characteristic manifesto problem: meaning often sits in the
relation between sections, not in either section alone.

This picture also explains why node labels matter. A full-document
expert score calibrates the root. A leaf or merge label checks
whether a smaller state still carries the same evidence for the
same dimension. Manifesto quasi-sentence codings are a natural fit
because they already mark local policy statements. They certify a
tree only when sampled or aggregated under a rule that targets the
same policy dimension as the root score.

\subsection{Benchmark, Pipeline, and References}
\label{sec:v8-benoit-benchmark}

The Manifesto Project collects party platforms, releases document
text through the Manifesto Corpus, and releases hand-coded
policy-emphasis data through the Manifesto Data Collection
\citep{ManifestoProject2025a,ManifestoCorpus2025}. Its coding
tradition goes back to the Manifesto Research Group and the
Comparative Manifestos Project
\citep{BudgeRobertsonHearl1987,BudgeEtAl2001,KlingemannEtAl2006,VolkensEtAl2013};
\citet{MerzEtAl2016} describe the modern corpus as a multilingual
text resource. The unit of human coding is the quasi-sentence, a
text unit expressing one political statement and assigned to a
policy category.

\citet{BenoitEtAl2025} introduce an LLM scoring pipeline against the
Manifesto Project expert means. Their pipeline first summarizes
each manifesto--dimension pair into a short English summary, then
scores that summary on the seven-point expert scale. We replicate
that pipeline with C-TreePO on the same task. We score a method by
the Pearson correlation between its root score and the expert mean
within each dimension; the macro score is the unweighted mean of
the six dimension correlations.

Three reference values from \citet{BenoitEtAl2025} anchor the
comparison: the split-expert agreement row, macro \(r=0.873\); the
proprietary LLM ensemble, macro \(r=0.817\); and the matched
Gemma-3 open-weight replication, macro \(r=0.793\). The split-expert
row is a measurement-agreement reference (experts are split into
two groups and the group means are correlated). We use it as
measurement context; no ladder is trained or selected on it.

\subsection{Tree Setup and Splits}\label{sec:v8-benoit-tree-setup}

The replication is a root-observed corpus evaluation. Each held-out
manifesto--dimension pair has a full-document expert target, and we
score the tree's root directly against that target. A tree run
chunks the manifesto into contiguous leaves of a fixed token size,
writes a short policy-evidence summary of each leaf with a learned
summarizer \(g\), merges adjacent summaries pairwise into parent
summaries with the same \(g\), continues up the tree until one root
summary represents the manifesto, and asks a learned scorer
\(\widehat f\) to read that root summary. We score the root against
expert means; a flat baseline asks the same scorer to read the
whole manifesto. The document-level target is unchanged from
\citet{BenoitEtAl2025}; what changes is the representation that the
scorer reads.

The leaf-size sweep runs the supervision-grain trade forward. Each
manifesto becomes a tree of policy-evidence units at multiple
grains: leaves of \(256\) to \(8{,}192\) tokens (roughly one
passage to one chapter), summaries of pairs of leaves, merges
further up, and the root. We report held-out Pearson at every
grain to expose how compression scales trade against the document
target, and the local diagnostic gap at each grain marks where the
proxy channel diverges from the held-out expert target.

We train \(\widehat f\) and \(g\) in alternating rounds: hold one
operator fixed, optimize the other against per-node training
targets, swap, repeat. The shorthand \(f^a g^b\) records the state
after \(a\) scorer-side and \(b\) summarizer-side updates. \(f^1
g^0\) is the first scorer trained against the initial summarizer;
\(f^1 g^1\) is the first summarizer trained against that scorer;
and so on. We report two points per leaf size below: the state
after round 1 and the best round at that leaf size.

Both \(\widehat f\) and \(g\) are DSPy programs that wrap the same
open-weight LLM with a prompt and a small bank of bootstrap
few-shot demonstrations.\footnote{Gemma-4-31B in NVFP4 (4-bit)
quantization. Re-running the pipeline on Gemma-3-27B-IT, the model
\citet{BenoitEtAl2025} use for their open-weight baseline, gives
Pearson correlations close to the numbers reported below; the gains
come from the tree structure, and not from a stronger underlying
model.} DSPy's MIPROv2 optimizer proposes and selects prompt
instructions and demonstrations against a per-node training metric
while the LLM weights stay frozen on the inference server. The
metric is mean per-node agreement against a teacher trace cached
before training: the initial scorer's score on the raw text spanned
by each node, computed once.
Section~\ref{sec:v8-benoit-local-diagnostics} unpacks the trace.

We tried four operator initializations during development. The
reported configuration starts both \(\widehat f^0\) and \(g^0\) as
default DSPy programs with rubric-style prompts and no optimized
demonstrations, and runs the alternation from there. Held-out
Pearson on the six dimensions stays in the same range across the
four starts at the leaf sizes we cover.
Appendix~\ref{app:v8-claim-tiers-manifesto} records the four init
designs in full.

The benchmark contains \(218\) manifestos with complete text,
expert targets, and teacher-trace preprocessing, split into \(140\)
training documents, \(30\) reserved documents (used for model
selection and diagnostics), and \(48\) held-out test documents.
Each manifesto is materialized as one tree per leaf size; reported
metrics are held-out root Pearson against expert-survey means. The
node count scales with leaf size: smaller leaves create more audit
and training units, larger leaves push more evidence into each
local state. Generated figures and tables live under
\texttt{assets/benoit/};
Appendix~\ref{app:v8-benoit-replication} records corpus, split,
benchmark, control, and artifact provenance.

\subsection{Single-Dimension Results}\label{sec:v8-benoit-single-dim}

Single-dimension training sets the per-dimension ceiling. We train
one \((\widehat f, g)\) pair per policy dimension against that
dimension's expert mean, sweep leaf size from \(256\) to \(8{,}192\)
tokens, and report two points per leaf: the state after round 1 and
the best round at that leaf size. The result is what the tree
representation can do on a dimension when one summarizer and one
scorer are dedicated to it, before we ask one shared summary to
serve all six dimensions at once.

\begin{figure}[!htbp]
  \centering
  \includegraphics[width=\linewidth]{manifesto_singledim_per_dim_live.pdf}
  \caption{\textbf{Single-dimension f/g results, all six policy dimensions.}
  External Pearson \(r\) against expert-survey means on the \(48\)-document
  held-out split, plotted against leaf token size. External means the score is
  computed on documents not used for training or model selection. Two series
  per panel: the state after round 1 and the best round at each leaf size.
  Reference lines are from \citet{BenoitEtAl2025}.}
  \label{fig:v8-benoit-singledim-per-dim}
\end{figure}

Figure~\ref{fig:v8-benoit-singledim-per-dim} resolves the per-dimension
ceiling. Five of the six dimensions clear the proprietary 18-score
ensemble in their best cells: economic policy (\(r=0.909\) at the
\(256\)-token leaf, after round 1), EU integration (\(r=0.965\)),
immigration (\(r=0.938\)), social policy (\(r=0.886\)), and
decentralization (\(r=0.557\) at the \(256\)-token leaf, ahead of the
ensemble's \(0.490\)). Environment sits within \(0.025\) of the
open-weight reference at its best leaf. The economic-policy ladder is
also the leaf-invariance test: held-out Pearson moves within a
\(0.04\)-wide range across \(32\times\) in leaf size, and the
\(256\)-token leaf clears split-expert agreement (\(0.880\)) by
\(0.029\).

Decentralization is the stress dimension. The evidence is less
consistent across parties and countries than economic-policy
evidence is. A party can endorse local administration while
opposing fiscal devolution. It can support regional identity while
keeping central budget authority. It can discuss federal reform in
country-specific institutional terms. The best round per leaf is
\(0.557\), \(0.476\), \(0.536\), \(0.479\), and \(0.485\) at
\(256\), \(512\), \(1{,}024\), \(2{,}048\), and \(4{,}096\) tokens.
Round-by-round the trajectory wanders: at \(256\) tokens, \(r\)
moves from \(0.534\) at \(f^1g^0\) up to \(0.557\); at \(512\)
tokens it slips from \(0.476\) to \(0.449\); at \(1{,}024\) tokens
it drops from \(0.510\) to \(0.435\) and recovers to \(0.536\) at
\(f^2g^2\). The \(256\)-token rung clears the proprietary 18-score
ensemble (\(0.490\)) by \(0.067\). The peak sits well below
split-expert agreement (\(0.780\)), and the local diagnostic gap
stays large at every leaf size. Decentralization is teacher-bounded:
the trained scorer fits its teacher trace better than the teacher
itself fits gold.

The dimension asks the state to retain institutional distinctions
that vary by country and by policy area: administrative devolution,
fiscal autonomy, regional identity, veto power, and central budget
authority. A scalar local score cannot carry those distinctions to
the parent. A readable state blurs them unless the state map is
trained on the dimension directly, and the teacher signal sets the
ceiling even then.

\subsection{Universal Summary}\label{sec:v8-benoit-universal-summary}

The universal-summary run probes the bounds of summary compression.
One state map \(g\) produces one summary per manifesto; six
per-dimension scorers all read it. One training run replaces six,
and the summary has to carry every scorer's evidence at once.

When the six dimensions ask about overlapping spans, that trade
pays off. One summary serves all of them. When they ask about
different spans, the summary has to compromise. A finite-bandwidth
string cannot represent six evidence geometries at full resolution,
and at least one scorer ends up reading a state that lacks the
evidence its dimension cares about. Section~\ref{sec:v8-ctree-math}
calls that a C1 failure in the representation. The training has a
separate cost: \(g\) is optimized against the average of the six
per-dimension rewards, so any one dimension's signal is diluted by
the other five, and the shared \(g\) moves slower than a dedicated
\(g\) would. Both costs show up below: on four of six dimensions
they are small, on decentralization they compound.

\(\widehat f\) stays dimension-aware. One shared training
trajectory produces all six per-dimension series across the leaf
sweep from \(256\) to \(8{,}192\) tokens. This is the application
closest to the C-TreePO objective: root expert anchors preserve the
document target, and teacher/proxy local-law losses push one state
map to retain evidence for multiple downstream readouts.

\begin{table}[t]
  \centering
  \resizebox{\linewidth}{!}{% Generated table; do not edit by hand.
% Metric: Pearson r (higher better)
\begin{tabular}{lrrrrrrrr}
\toprule
Method & Econ. & Social & Immig. & EU & Env. & Decent. & Macro & \#/6 \\
\midrule
\multicolumn{9}{l}{\textit{Benoit et al. (2025) reference}} \\
Proprietary ensemble, 18 scores (Fig 1) & $+0.870$ & $+0.920$ & $+0.890$ & $+0.910$ & $+0.820$ & $+0.490$ & $+0.817$ & 6/6 \\
Split-expert reference (Table 3) & $+0.880$ & $+0.910$ & $+0.880$ & $+0.950$ & $+0.840$ & $+0.780$ & $+0.873$ & 6/6 \\
LLaMA-3.3-70B (Table 6) & $+0.840$ & $+0.870$ & $+0.860$ & $+0.860$ & $+0.680$ & $+0.400$ & $+0.752$ & 6/6 \\
DeepSeek-V3 (Table 6) & $+0.840$ & $+0.870$ & $+0.890$ & $+0.860$ & $+0.790$ & $+0.450$ & $+0.783$ & 6/6 \\
Gemma-3-27B-IT (Table 6) & $+0.860$ & $+0.860$ & $+0.890$ & $+0.840$ & $+0.860$ & $+0.450$ & $+0.793$ & 6/6 \\
\midrule
\multicolumn{9}{l}{\textit{Ours: per-dim pipeline $\times$ leaf size (Gemma-4-31B-NVFP4, 1 summarizer + 1 scorer per dim)}} \\
leaf = 64 K chars ($\approx$16K tokens) & $+0.916$ & $+0.817$ & $+0.888$ & $+0.924$ & $+0.854$ & $+0.489$ & $+0.815$ & 6/6 \\
leaf = 32 K chars ($\approx$8K tokens) & $+0.929$ & $+0.857$ & $+0.886$ & $+0.916$ & $+0.877$ & $+0.480$ & $+0.824$ & 6/6 \\
leaf = 24 K chars ($\approx$6K tokens) --- full test n$\approx$215 & $+0.892$ & $+0.840$ & $+0.867$ & $+0.896$ & $+0.814$ & $+0.464$ & $+0.796$ & 6/6 \\
leaf = 16 K chars ($\approx$4K tokens) & $+0.925$ & $+0.847$ & $+0.898$ & $+0.909$ & $+0.831$ & $+0.537$ & $+0.824$ & 6/6 \\
leaf =  8 K chars ($\approx$2K tokens) & $+0.939$ & $+0.869$ & $+0.884$ & $+0.901$ & $+0.839$ & $+0.543$ & $+0.829$ & 6/6 \\
\midrule
\multicolumn{9}{l}{\textit{Ours: combined pipeline $\times$ leaf size (one shared summarizer with union rubric, 6 scores)}} \\
leaf = 64 K chars & $+0.910$ & $+0.867$ & $+0.897$ & $+0.910$ & $+0.795$ & $+0.458$ & $+0.806$ & 6/6 \\
leaf = 32 K chars & $+0.917$ & $+0.849$ & $+0.908$ & $+0.889$ & $+0.749$ & $+0.438$ & $+0.792$ & 6/6 \\
leaf = 24 K chars (full test n=229) & $+0.868$ & $+0.850$ & $+0.880$ & $+0.902$ & $+0.809$ & $+0.396$ & $+0.784$ & 6/6 \\
leaf = 16 K chars & $+0.919$ & $+0.851$ & $+0.916$ & $+0.889$ & $+0.808$ & $+0.441$ & $+0.804$ & 6/6 \\
leaf =  8 K chars & $+0.927$ & $+0.874$ & $+0.911$ & $+0.917$ & $+0.801$ & $+0.413$ & $+0.807$ & 6/6 \\
\midrule
\multicolumn{9}{l}{\textit{Ours: tiny-leaf extensions of the chunk sweep (Gemma-4)}} \\
tree, leaf = 4 K chars & $+0.932$ & $+0.886$ & $+0.885$ & $+0.931$ & $+0.838$ & $+0.580$ & $+0.842$ & 6/6 \\
tree, leaf = 2 K chars & $+0.924$ & $+0.867$ & $+0.887$ & $+0.906$ & $+0.793$ & $+0.584$ & $+0.827$ & 6/6 \\
tree, leaf = 1 K chars (stress test, depth $\geq$6) & $+0.933$ & $+0.859$ & $+0.914$ & $+0.925$ & $+0.838$ & $+0.493$ & $+0.827$ & 6/6 \\
\midrule
\multicolumn{9}{l}{\textit{Ours: concat-no-merge (chunks summarized independently, joined, scored --- tests whether the merge step carries signal)}} \\
concat, leaf = 32 K chars & $+0.929$ & $+0.856$ & $+0.897$ & $+0.920$ & $+0.848$ & $+0.571$ & $+0.837$ & 6/6 \\
concat, leaf = 16 K chars (default) & $+0.931$ & $+0.823$ & $+0.909$ & $+0.912$ & $+0.866$ & $+0.523$ & $+0.827$ & 6/6 \\
concat, leaf =  8 K chars & $+0.932$ & $+0.849$ & $+0.882$ & $+0.915$ & $+0.850$ & $+0.594$ & $+0.837$ & 6/6 \\
\midrule
\multicolumn{9}{l}{\textit{Ours: flat baseline (no chunk, no summary; truncate text and score) --- tests whether tree is needed at all}} \\
flat, truncation = 48 K chars & $+0.936$ & $+0.821$ & $+0.926$ & $+0.889$ & $+0.806$ & $+0.587$ & $+0.827$ & 6/6 \\
flat, truncation = 24 K chars (default) & $+0.929$ & $+0.847$ & $+0.889$ & $+0.882$ & $+0.829$ & $+0.542$ & $+0.820$ & 6/6 \\
flat, truncation = 12 K chars & $+0.929$ & $+0.844$ & $+0.938$ & $+0.880$ & $+0.779$ & $+0.564$ & $+0.822$ & 6/6 \\
flat, truncation =  6 K chars & $+0.926$ & $+0.811$ & $+0.967$ & $+0.881$ & $+0.886$ & $+0.243$ & $+0.786$ & 6/6 \\
\midrule
\multicolumn{9}{l}{\textit{Ours: scorer ablations (Gemma-4, Benoit's GPT-4o summaries held fixed)}} \\
1. Zero-shot scorer & $+0.832$ & $+0.886$ & $+0.858$ & $+0.905$ & $+0.668$ & $+0.311$ & $+0.743$ & 6/6 \\
2. Few-shot optimized scorer & $+0.822$ & $+0.892$ & $+0.853$ & $+0.912$ & $+0.627$ & $+0.297$ & $+0.734$ & 6/6 \\
3. Joint scorer baseline (shared across 6 dims) & $+0.837$ & $+0.881$ & $+0.852$ & $+0.904$ & $+0.658$ & $+0.314$ & $+0.741$ & 6/6 \\
4. Joint scorer, few-shot optimized & $+0.817$ & $+0.891$ & $+0.848$ & $+0.908$ & $+0.700$ & $+0.313$ & $+0.746$ & 6/6 \\
5. Joint scorer, reflective prompt-optimized & $+0.832$ & $+0.882$ & $+0.848$ & $+0.879$ & $+0.688$ & $+0.344$ & $+0.746$ & 6/6 \\
\midrule
\multicolumn{9}{l}{\textit{Ours: exact-model replication (Benoit's Gemma-3-27B-IT BF16)}} \\
Gemma-3-27B scorer-only, OUR scoring rubric & $+0.826$ & $+0.864$ & $+0.846$ & $+0.902$ & $+0.573$ & $+0.329$ & $+0.723$ & 6/6 \\
Gemma-3-27B scorer-only, exact Benoit rubric & $+0.814$ & $+0.895$ & $+0.852$ & $+0.852$ & $+0.616$ & $+0.359$ & $+0.731$ & 6/6 \\
Gemma-3-27B, raw-prompt Benoit format (bare-integer response) & $+0.782$ & $+0.901$ & $+0.854$ & $+0.874$ & $+0.590$ & $+0.368$ & $+0.728$ & 6/6 \\
\bottomrule
\end{tabular}
}
  \caption{\textbf{Pearson \(r\) against the expert-survey means.}
  Reference rows reproduce the \citet{BenoitEtAl2025} pipeline; C-Tree
  rows use the learned manifesto compression pipeline.}
  \label{tab:v8-benoit-headline}
\end{table}

The per-dimension tree at \(8\)K-character leaves reaches macro
\(r=0.829\) (Table~\ref{tab:v8-benoit-headline}), with dimension
scores \(0.939\) economic, \(0.869\) social, \(0.884\) immigration,
\(0.901\) European integration, \(0.839\) environment, and \(0.543\)
decentralization. The universal-summary tree at the same leaf size
reaches macro \(r=0.807\), with \(0.927\), \(0.874\), \(0.911\),
\(0.917\), \(0.801\), and \(0.413\) on the same six dimensions. The
macro gap of \(0.022\) hides the dimension-level structure: economic,
social, immigration, and EU integration sit within \(0.030\) of their
per-dimension peaks under the universal summary; decentralization
sits \(0.130\) below.

\begin{figure}[!htbp]
  \centering
  \includegraphics[width=\linewidth]{manifesto_fg_combined_ladder_f1g0_f1g1.pdf}
  \caption{\textbf{Universal-summary run, per-dimension trajectories.}
  External Pearson \(r\) against expert-survey means on the \(48\)-document
  held-out split from one shared \((\widehat f,g)\) training trajectory. Two
  series per panel: the state after round 1 and the best round at each leaf
  size.}
  \label{fig:v8-benoit-combined-ladder}
\end{figure}

\begin{table}[t]
  \centering
  \begin{tabular}{lrrrr}
    \toprule
    Dimension & Best \(r\) & Leaf tokens & Stage & \(n\) \\
    \midrule
    Economic & \(0.886\) & \(4{,}096\) & \(f^2g^1\) & \(47\) \\
    Social & \(0.886\) & \(1{,}024\) & \(f^2g^2\) & \(43\) \\
    Immigration & \(0.938\) & \(4{,}096\) & \(f^3g^3\) & \(29\) \\
    European integration & \(0.965\) & \(2{,}048\) & \(f^1g^1\) & \(35\) \\
    Environment & \(0.795\) & \(1{,}024\) & \(f^1g^1\) & \(42\) \\
    Decentralization & \(0.461\) & \(8{,}096\) & \(f^1g^1\) & \(42\) \\
    \bottomrule
  \end{tabular}
  \caption{\textbf{Best per-dimension cells under the universal summary.}
  Leaf tokens is the token budget for one leaf, stage is the f/g ladder state,
  and \(n\) is the number of held-out documents with usable labels for that
  dimension. A common deployed row would choose one leaf size and one stage;
  this table shows the best cell per dimension to diagnose state-sharing
  limits.}
  \label{tab:v8-benoit-joint-dim-best}
\end{table}

Figure~\ref{fig:v8-benoit-combined-ladder} and
Table~\ref{tab:v8-benoit-joint-dim-best} read the same pattern at the
trajectory level. Where a single-dimension run exists for comparison,
the result is sharp. Economic's universal-summary peak (\(r=0.886\) at
leaf \(4{,}096\)) sits within \(0.002\) of its single-dimension peak
(\(0.888\) at leaf \(2{,}048\)); the universal summary captures
essentially the same signal as a dedicated economic summary.
Decentralization's universal-summary peak (\(r=0.461\) at leaf
\(8{,}192\)) sits \(0.096\) below its single-dimension peak (\(0.557\)
at leaf \(256\)); the single-dimension peak also clears the proprietary
ensemble on this dimension by \(0.067\).

The universal-summary run shows the same decentralization shape at
every leaf size. Held-out Pearson stays between \(0.30\) and
\(0.39\), the local diagnostic gap stays around \(0.45\), and the
stage-by-stage trajectory oscillates: a \(g\)-update lifts the score
and the next \(\widehat f\)-update pulls it back. The training cost
shows up as multi-task interference: averaging the six scorer-side
rewards leaves any single dimension with a fraction of the gradient
when its target points away from the others, so a \(g\)-step that
picks up that dimension's signal gets unwound by the next
\(\widehat f\)-step. Dropping the averaging (the single-dimension
run, Section~\ref{sec:v8-benoit-single-dim}) lifts held-out Pearson
on decentralization from \(0.461\) to \(0.557\) and recovers most of
the gap. The representation cost is the structural piece. A
universal summary trained to serve all six dimensions can preserve
the shared evidence well and still blur the
decentralization-specific distinctions; a dimension whose evidence
diverges that far from the shared average needs more capacity,
different law weights, or its own state map.

Environment is a different story. A cold-start control makes the
source visible: scoring stored full-document environment summaries
with the open-weight LLM and the raw \citet{BenoitEtAl2025} prompt
reaches \(r=0.784\) on the held-out split, essentially tied with
the first ladder row at \(256\) tokens. The ceiling sits with the
underlying scorer and the prompt; the universal-summary ladder
preserves the cold-start signal and cannot manufacture more.

\subsection{Local Diagnostics and Node-Label Route}
\label{sec:v8-benoit-local-diagnostics}

The local diagnostic gap is the audit-design tool in practice. It
identifies where proxy supervision diverges from the held-out
expert target, which is where sampled \(f^\ast\) calls would
correct the most error. The gap is computed from teacher-trace
residuals: at each node, the discrepancy between a recorded teacher
local target and the current proxy/readout behavior. The teacher
trace itself is the initial scorer's prediction on the raw text
spanned by each node, cached once before training; agreement
against this trace is what the training metric in
Section~\ref{sec:v8-benoit-tree-setup} maximizes. The proxy
\(\widehat f\) produces local-law losses everywhere; sampled
\(f^\ast\) node calls then correct the proxy channel through
Equation~\eqref{eq:v8-corrected-node-loss}.

\begin{figure}[!htbp]
  \centering
  \includegraphics[width=\linewidth]{manifesto_fg_combined_audit_gap.pdf}
  \caption{\textbf{Universal-summary local diagnostic gaps.}
  Per-dimension teacher-trace gaps from the universal-summary run.
  Decentralization's gap is the main warning signal.}
  \label{fig:v8-benoit-combined-audit-gap}
\end{figure}

\begin{figure}[!htbp]
  \centering
  \includegraphics[width=\linewidth]{manifesto_singledim_per_dim_live_audit_gap.pdf}
  \caption{\textbf{Single-dimension local diagnostic gaps.}
  Teacher-trace gap for the single-dimension ladders in
  Figure~\ref{fig:v8-benoit-singledim-per-dim}.}
  \label{fig:v8-benoit-single-audit-gap}
\end{figure}

The diagnostic flagged decentralization's failure mode on its own,
without the single-dimension comparison being available. The proxy
tracks its teacher trace tightly while the trained scorer tracks
gold loosely, and the gap is widest where shared compression
averages over dimensions whose evidence geometry differs
(Figure~\ref{fig:v8-benoit-combined-audit-gap}). That is the
diagnostic and the corrected-node-loss working as a pair: the gap
points at the nodes where sampled \(f^\ast\) calls would estimate
and correct \(b_v=\widehat\ell_v-\ell_v^\ast\) at the highest
return per call.

The reported Manifesto numbers are root-observed scores plus
teacher-trace local diagnostics. The node-audit path aggregates
Manifesto Project quasi-sentence evidence to each sampled span
under a declared dimension-specific rule, evaluates C1/C2/C3 losses
on those sampled nodes, and applies the logged-propensity
correction from Equation~\eqref{eq:v8-corrected-node-loss}.
Appendix~\ref{app:v8-manifesto-audit-design} specifies that design.

\subsection{Where It Works, What It Doesn't, and Where Leaves Help}
\label{sec:v8-benoit-takeaway}

The reported run sorts cleanly. Per-dimension trees reach macro
\(r=0.829\); the universal-summary tree reaches \(r=0.807\). Both
clear the matched open-weight baseline (\(0.793\)); the per-dimension
tree also clears the proprietary 18-score ensemble (\(0.817\)) at
the root. A direct full-document control on the same scorer LLM
(Gemma-4 reading the entire manifesto under a MIPRO-tuned prompt)
reaches macro \(r=0.841\), with the per-dimension breakdown
\(0.968\) economic, \(0.884\) social, \(0.631\) decentralization,
\(0.874\) environment, \(0.779\) EU, \(0.912\) immigration. The
tree-compressed root reads sit within \(0.012\) of that
full-document ceiling on the macro score. On the five well-behaved
dimensions, held-out Pearson under the universal summary stays
within \(0.030\) of the per-dimension peak at every leaf size,
and the leaf-size sweep moves the per-dimension trees within a
\(0.04\)-wide band on economic policy across \(32\times\) in leaf
size. An analyst can spend the same audit budget on shorter, more
numerous units or on longer, fewer ones and reach the same
document-level target.

Decentralization is the dimension that doesn't sort the same way.
Held-out Pearson stays well below split-expert agreement at every
leaf size, the trajectory wanders round-by-round, and the
universal-summary cost over the per-dimension peak is large.

Where leaves help is visible in the same run. The local diagnostic
gap is widest exactly where shared compression averages over
dimensions whose evidence geometry differs, and it flagged the
decentralization breakdown without the single-dimension comparison
being available. The gap and the corrected-node-loss work as a
pair: the gap points at the nodes where sampled \(f^\ast\) calls
would estimate and correct \(b_v=\widehat\ell_v-\ell_v^\ast\) at
the highest return per call. A node audit is the next labeling
decision, not another root score.

Three failure modes account for the cases where the framework
breaks. The first is a disagreement between root accuracy and
local faithfulness. A tree can predict the right root score while
losing the evidence an expert would want to inspect. That matters
for measurement, because the same score is then used in downstream
inference, in cross-country comparisons, and in arguments about
which claims drove the placement. Those uses are inspectable only
when the local states remain faithful.

The second is state sharing across incompatible dimensions. The
universal-summary run asks one policy state to support six
readouts. That is efficient when dimensions share evidence, but it
dilutes a dimension whose evidence geometry is different.
Decentralization is the warning case because it depends on
institutional distinctions, country-specific arrangements, and
boundary-spanning qualifications that are easy to compress away.
The available fixes are a dimension-specific state map, a higher
local-law weight for that dimension, or a sampling design that
oversamples its high-risk nodes; a larger backbone model is not
the obvious first step.

The third is proxy overconfidence. A dense LLM judge can provide
smooth local losses everywhere, but if its errors are correlated
with the same distinctions that experts care about, the local
channel can train the state map toward the wrong geometry. On
economic policy, this might look like an LLM judge that
systematically under-weights labor-market language and over-weights
tax language; the local channel then trains the state map toward
tax-emphasis evidence at the cost of labor evidence, even when the
root score still tracks the expert mean. The audit gap is what
makes that drift visible.
Equation~\eqref{eq:v8-corrected-node-loss} uses sampled \(f^\ast\)
nodes to estimate and correct exactly this proxy bias.

\subsection{When Local Labels Pay}
\label{sec:v8-benoit-when-locals-pay}

Local labels pay when three conditions hold. First, the target has
meaningful local evidence. Manifesto policy scores usually do: tax
pledges, welfare commitments, immigration restrictions, European
integration claims, and regional-autonomy statements appear in
identifiable spans. Second, the tree state has enough capacity to
carry cross-span interactions, such as a tax cut qualified by a
later welfare promise or a devolution claim limited by central
budget authority. Third, the local labels measure the same
construct as the root target.

Under those conditions, node labels are more than cheaper labels.
They tell the learner what to preserve before the root score is
read. They also tell the auditor where the root score is fragile.
If changing leaf size leaves the root estimate stable, smaller
leaves create more auditable sites without changing the estimand.
If a local diagnostic gap is concentrated in one dimension, as it
is for decentralization here, the next label should usually go to
that dimension's nodes rather than to another random root score.

Local labels pay less when calibration, not sampling, is the
bottleneck. If the accessible local target is misaligned with the
expert-survey dimension, adding more of it shrinks
\(B_{\mathrm{est}}\) while leaving \(B_{\mathrm{cal}}\) large.
They also pay less when the target is intrinsically global in a
way the state map cannot localize. A rhetorical-style judgment or
a party-brand interpretation may depend on document-wide framing
that no span-level evidence can preserve. In that case the audit
should show high C1 or C3 distortion, and the certificate should
widen rather than manufacture confidence.
